% ---------------------------------------------------------------------------
%  sec_howitworks.tex -- how the machine works: the bricks, the unknowns and
%  the curl matrix. Sits between the aims and the formal FIT section.
% ---------------------------------------------------------------------------
\section{How the machine works}\label{sec:howitworks}

This section is the tour we would give a colleague who wants to know what the solver does before reading the
formal statement of it. It describes the geometry the program sees, the numbers it stores on that geometry, and
the single matrix that turns a statement about one loop into a statement about every loop at once. There is no
derivation here and nothing is measured: the derivation is in Section~\ref{sec:theory}, and the evidence that
back the claims below is in the section on the discrete energy identity and in the section on the cavity
resonances. Nothing in this section is approximate except one thing, which we will point at when we come to it,
and that thing is the material. The topology -- what touches what, which edge bounds which face, what a loop
integral means -- is either right or wrong, and in this discretisation it is right.

\subsection{A box, cut into bricks}\label{sec:hiw:bricks}

The computational domain is a box, and the first decision, the one everything else follows from, is to cut it
into bricks. A brick is $\Delta x$ by $\Delta y$ by $\Delta z$; the three spacings may differ, and they may
differ from one brick to the next, but every brick is a box and every internal brick face is shared, whole, with
the face of exactly one neighbour. There is no other geometry in the program. The solver does not know that the
thing inside the box is a cavity with a fillet in one corner or a patch of laminate; it knows the bricks, the
numbering of their edges and faces, and which faces have been marked as conductor. When we say that a structure
has been discretised, this is the entirety of what we mean.

A single brick has eight corners, twelve edges and six faces. Interior faces are shared by two bricks and
interior edges by four, and it is that sharing rather than the brick itself which fixes the counts. For a mesh
of $N$ bricks there are, to within boundary terms that do not change the picture, $N$ corners, $3N$ edges and
$3N$ faces. The corners are the vertices, the edges join them, the faces are spanned by the edges and the bricks
are the cells the faces enclose; Figure~\ref{fig:hiw:brick} draws one brick in this reading.

Two of those counts are equal, and that is the first thing worth pausing on. A mesh with $3N$ edges has $3N$
faces, not approximately and not only when the bricks are cubes. It is the discrete form of the exact sequence
that runs vertices to edges to faces to cells in the continuum, where the boundary of a face is a loop of edges
and the boundary of a cell is a closed surface of faces, and where there is nothing left over at either end.
The practical consequence is that there are no spare slots and no choice about where the field goes. It has to
live on the edges and on the faces, because those are the entities whose counts match; an electric field held
at the corners would be short of unknowns by a factor of three in every direction, and there would be no loop
for Faraday's law to be written on.

What the solver knows of geometry is therefore the mesh and nothing else. A perfect electric conductor is a set
of faces marked as such, and the four edges bounding each of those faces are known quantities rather than
extrapolated ones. There is no analytic surface anywhere in the code to be intersected or sampled, which is why
the same code path takes a mesh with unequal spacings in all three directions without any change of
formulation. The shape of the object enters once, through the mesh; given the mesh, everything below is exact.

\begin{figure}[htbp]
\centering
\begin{tikzpicture}[font=\footnotesize]
  % the shaded face: the top of the brick
  \draw[face] (0,2.2) -- (3.6,2.2) -- (5.1,3.3) -- (1.5,3.3) -- cycle;
  % the three edges hidden behind the body, left thin
  \draw[mesh thin] (0,0) -- (1.5,1.1);
  \draw[mesh thin] (1.5,1.1) -- (5.1,1.1);
  \draw[mesh thin] (1.5,1.1) -- (1.5,3.3);
  % the nine visible edges
  \draw[primal] (0,0) -- (3.6,0) -- (3.6,2.2) -- (0,2.2) -- cycle;
  \draw[primal] (3.6,0) -- (5.1,1.1) -- (5.1,3.3);
  \draw[primal] (3.6,2.2) -- (5.1,3.3);
  \draw[primal] (0,2.2) -- (1.5,3.3) -- (5.1,3.3);
  % one edge unknown
  \draw[edof] (0.35,0) -- (3.25,0);
  \node[lbl, anchor=south] at (1.8,0.12) {$e=\int_L \vec E\cdot \mathrm{d}\vec l\ (\si{\volt})$};
  % one face unknown and the direction of its normal
  \draw[harrow] (2.55,2.75) -- (2.55,3.78);
  \node[lbl, anchor=west] at (2.65,3.72) {$b=\iint_A \vec B\cdot \mathrm{d}\vec A\ (\si{\weber})$};
  % dimensions
  \draw[dimline] (0,-0.7) -- (3.6,-0.7);
  \node[lbl, anchor=north] at (1.8,-0.8) {$\Delta x$};
  \draw[dimline] (-0.7,0) -- (-0.7,2.2);
  \node[lbl, anchor=east] at (-0.8,1.1) {$\Delta y$};
  \draw[dimline] (4.1,-0.75) -- (5.65,0.4);
  \node[lbl, anchor=west] at (5.75,0.45) {$\Delta z$};
\end{tikzpicture}
\caption{One brick of a structured mesh, drawn in cabinet projection with the three hidden edges left thin so
that all twelve are accounted for. The unknown attached to an edge is the line integral of the electric field
along the whole edge, marked in red, and the unknown attached to a face is the flux of the magnetic field
through that face, marked in blue on the shaded top face. Neither mark is a point sample: the red arrow stands
for a definite integral along an edge, and the blue arrow gives the direction of a surface normal, nothing
more.}
\label{fig:hiw:brick}
\end{figure}

\subsection{What lives on an edge, and what lives on a face}\label{sec:hiw:dofs}

The unknown we store on an edge is not the electric field at some point of that edge. It is the line integral of
the electric field along the edge,
\begin{equation}
  e_i = \int_{L_i} \vec E \cdot \mathrm{d}\vec l ,
  \label{eq:hiw-e}
\end{equation}
in volts, and the unknown we store on a face is the surface integral of the magnetic flux density through it,
\begin{equation}
  b_j = \iint_{A_j} \vec B \cdot \mathrm{d}\vec A ,
  \label{eq:hiw-b}
\end{equation}
in webers. These two choices are the heart of the method, and they are not a convenience adopted to make the
algebra tidy.

Faraday's law is a statement about a loop and the surface the loop bounds; Ampere's law is a statement about a
loop and the surface that pierces it. If the unknowns are the integrals those laws are written in terms of,
then each law becomes a relation between unknowns with nothing in between: no quadrature, no interpolation, no
Taylor expansion, no assumption that the field is smooth across a brick. Note what the first equation is not. It
is not a point value of $E$ multiplied by a length; $E$ may vary arbitrarily along the edge and the integral
still means exactly what it says. The single number $e_i$ is all the solver ever knows about the field on that
edge, and on a given mesh it is all the solver needs.

There is a practical argument for the same choice. A voltmeter returns a line integral, and a loop of wire
returns the rate of change of a flux through a surface, so the quantities a measurement hands us are the
quantities stored here. A point value would be an interpolant, manufactured from the integrals whenever a plot
wanted one; the solver itself never needs it, because it goes from edge integrals to face integrals and back to
edge integrals without ever evaluating $E$ at a point. When we do draw a field map, the interpolation happens in
the plotting code, where it can be seen for what it is.

The two sets of unknowns do not share a home. The $e$ unknowns live on the edges, which form one lattice; the
$b$ unknowns live at the face centres, which form a second lattice displaced by half a brick in each of the
three directions. In two dimensions this is exactly the staggered grid of elementary finite-difference practice,
and Figure~\ref{fig:hiw:staggered} draws the two lattices on the same patch. Nothing is collocated: no $e$ and
no $b$ occupy the same location, and no two unknowns of any kind sit at the same point. This is a feature and
not an accident of implementation. The offset is what allows a difference of two neighbouring values to stand
for a derivative without ever pretending that two quantities defined at the same place can differ.

The offset also makes boundary conditions answerable. A perfect electric conductor is a set of marked faces, the
tangential electric field vanishes on those faces, and in these unknowns that is four numbers set to zero -- a
statement about a loop which is true by construction and needs no extrapolation from the interior.

\begin{figure}[htbp]
\centering
\begin{tikzpicture}[font=\footnotesize]
  \def\s{0.8}
  % (a) the primal lattice and its edge unknowns
  \foreach \i in {0,...,4}{
    \draw[primal] ({\i*\s},0) -- ({\i*\s},{4*\s});
    \draw[primal] (0,{\i*\s}) -- ({4*\s},{\i*\s});
  }
  \foreach \j in {0,...,4}{
    \foreach \i in {0,...,3}{
      \node[edot] at ({(\i+0.5)*\s},{\j*\s}) {};
      \node[edot] at ({\j*\s},{(\i+0.5)*\s}) {};
    }
  }
  \draw[edof] ({1.1*\s},{1*\s}) -- ({1.9*\s},{1*\s});
  \draw[edof] ({2*\s},{1.1*\s}) -- ({2*\s},{1.9*\s});
  \draw[edof] ({1.9*\s},{2*\s}) -- ({1.1*\s},{2*\s});
  \draw[edof] ({1*\s},{1.9*\s}) -- ({1*\s},{1.1*\s});
  \node[lbl, anchor=north] at ({2*\s},-0.18) {(a) primal lattice: $e$ on every edge};
  % separator
  \draw[mesh thin] (4.0,-0.35) -- (4.0,{4*\s+0.05});
  % (b) the dual lattice, shifted by half a cell
  \begin{scope}[xshift=4.8cm]
    \foreach \i in {0,...,3}{
      \draw[dual] ({(\i+0.5)*\s},{0.5*\s}) -- ({(\i+0.5)*\s},{3.5*\s});
      \draw[dual] ({0.5*\s},{(\i+0.5)*\s}) -- ({3.5*\s},{(\i+0.5)*\s});
    }
    \foreach \j in {0,...,3}{
      \foreach \i in {0,...,3}{
        \node[bdot] at ({(\i+0.5)*\s},{(\j+0.5)*\s}) {};
      }
    }
    \draw[dual] ({1.5*\s},{2.5*\s}) circle (3.6pt);
    \node[bdot] at ({1.5*\s},{2.5*\s}) {};
    \node[lbl, anchor=north] at ({2*\s},-0.18) {(b) dual lattice: $b$ on every face};
    \node[lbl, anchor=north] at ({2*\s},-0.62) {ringed dot: normal out of the page};
  \end{scope}
\end{tikzpicture}
\caption{The two lattices in two dimensions, on the same four-by-four patch and at the same scale. In (a) the
primal lattice carries the field on every edge: the arrows are drawn on the edges of one cell and the dots sit
at the midpoints of all of them. In (b) the dual lattice, displaced by half a cell, carries the flux at the
centre of every face, one of which is ringed to show its normal pointing out of the page. The thing to look at
is that no red dot and no blue dot coincides.}
\label{fig:hiw:staggered}
\end{figure}

\subsection{The matrix whose entries are only ever plus one, minus one, or zero}\label{sec:hiw:curl}

Take one face of one brick and the four edges that bound it. Faraday's law for that face, written in the
integrals we have chosen, reads
\begin{equation}
  e_1 + e_2 + e_3 + e_4 = -\frac{\mathrm{d}}{\mathrm{d}t} \, b ,
  \label{eq:hiw-faraday-face}
\end{equation}
where the left-hand side is the circulation of $\vec E$ round the loop provided each $e_i$ is taken in the sense
in which the loop is traversed. Each edge of the mesh has a positive direction fixed once and for all by the
grid axes, so an edge whose stored direction agrees with the traversal contributes $+e_i$ and an edge that runs
against it contributes $-e_i$. That is the entire origin of the signs in this method. A sign is either right or
wrong; there is no third possibility for it to approximate.

Now do it for every face at once. Collect the edge unknowns into a column vector $\mathbf{e}$, one entry per
edge, and the face unknowns into a column vector $\mathbf{b}$, one entry per face. Each of the equations above
is a sum of a few entries of $\mathbf{e}$ with signs, and a sum of that kind is one row of a matrix multiplying
$\mathbf{e}$. The matrix is $\Cmat$: one row per face, one column per edge, with
\begin{equation}
  \Cmat_{ji} =
  \begin{cases}
    +1 & \text{edge $i$ bounds face $j$ and runs with its orientation,}\\
    -1 & \text{edge $i$ bounds face $j$ and runs against it,}\\
     0 & \text{edge $i$ does not bound face $j$,}
  \end{cases}
  \label{eq:hiw-C}
\end{equation}
and the whole of Faraday's law, for every face simultaneously, is the single equation
\begin{equation}
  \Cmat \mathbf{e} = -\frac{\mathrm{d}}{\mathrm{d}t}\,\mathbf{b} .
  \label{eq:hiw-faraday}
\end{equation}
Nothing has been truncated on the way to this line. It is the discrete statement, not a discretisation of one.

Ampere's law uses the same matrix read the other way,
\begin{equation}
  \Cmat^{\!\top} \mathbf{h} = \frac{\mathrm{d}}{\mathrm{d}t}\,\mathbf{d} + \mathbf{j} ,
  \label{eq:hiw-ampere}
\end{equation}
with $\mathbf{h}$ on the dual edges that pierce the primal faces, $\mathbf{d}$ on the dual faces that cross the
primal edges, and $\mathbf{j}$ the current through those dual faces. The transpose is not a second operator that
happens to look similar. It is the same incidence counted from the other side, because the two complexes
interlace: a dual edge pierces a primal face, and the four dual faces that cross the primal edges bounding that
face are exactly the dual faces which that dual edge bounds. The set of incidences is one set, and the entries
of $\Cmat^{\top}$ are the entries of $\Cmat$ -- still $\pm 1$, still zero.

Everything now follows from the sparsity pattern, and this is where the method earns its keep. Write $\Gmat$ for
the matrix that takes values at the vertices to differences along the edges; its entries are also only $\pm 1$
or $0$. Go round any closed loop of edges: each vertex is entered once and left once, so each vertex value is
added once and subtracted once, and the sum telescopes to nothing. Therefore
\begin{equation}
  \Cmat \Gmat = \mathbf{0} \qquad \text{and, one complex over,} \qquad \Smat \Cmat = \mathbf{0} ,
  \label{eq:hiw-identities}
\end{equation}
identically, for any mesh, with no remainder and no error term. These are the discrete curl-grad and div-curl
identities, and they are properties of the matrices rather than statements that become truer as the mesh is
refined. \emph{The topology is exact; the only approximation is the material.} The material is where the
approximation lives because the constitutive relations are the only place where a continuum has to be replaced
by a number: the entries of the material matrices are averages over a brick, and an average is a judgement about
a shape that the mesh resolves only coarsely.

The consequences are not separate results, they are that same sentence said three ways. The discrete energy
identity holds term by term rather than to within a truncation error, because the work done by the curl term
telescopes across the mesh exactly; in a lossless cavity the measured relative drift of the discrete energy is
at the level of $10^{-15}$, which is round-off and nothing else, as reported in Section~\ref{sec:energy}. The
cavity eigenvalues converge at second order, with a measured rate of $0.507$ against a predicted $0.510$ in
Section~\ref{sec:eigen}. We refer to those sections rather than repeat them, because the point being made here
is the shape of the argument: exact algebra, an error that can only enter through the material numbers, and a
convergence rate that follows.

It is fair to ask what all this buys, since on a Cartesian mesh the arithmetic that comes out is the arithmetic
Yee wrote down in 1966 -- the same neighbours, the same stencil, later the same leapfrog update. If a stencil
were all that was wanted, the finite-difference reading would have got there first and more cheaply. What the
integral derivation buys is meaning. Each unknown is a definite physical quantity with a unit attached, so a
boundary condition is a statement about a surface or a loop rather than about the first row of a stencil, and
averaging a material over a brick is a well-posed question with a right answer rather than an accident of which
point the stencil happened to sample. And the algebra survives a change of mesh shape: the same incidence
matrices, with the same identities, describe distorted or non-orthogonal cells, and only the material matrices
change. The finite-difference reading has no such continuation, because its stencil was written for a Cartesian
grid in the first place.

\begin{figure}[htbp]
\centering
\begin{minipage}[c]{0.50\linewidth}
\centering
\begin{tikzpicture}[font=\footnotesize]
  \draw[face] (0,0) -- (3.0,0) -- (3.6,2.0) -- (0.6,2.0) -- cycle;
  \draw[primal] (0,0) -- (3.0,0) -- (3.6,2.0) -- (0.6,2.0) -- cycle;
  \draw[edof] (0.2,0) -- (2.8,0);
  \draw[edof] (3.057,0.192) -- (3.543,1.808);
  \draw[edof] (3.4,2.0) -- (0.8,2.0);
  \draw[edof] (0.543,1.808) -- (0.057,0.192);
  \node[lbl, anchor=north] at (1.5,-0.12) {$e_1$};
  \node[lbl, anchor=west]  at (3.68,1.0) {$e_2$};
  \node[lbl, anchor=south] at (2.1,2.12) {$e_3$};
  \node[lbl, anchor=east]  at (0.02,1.0) {$e_4$};
  \draw[dual] (1.8,1.0) circle (3.6pt);
  \node[bdot] at (1.8,1.0) {};
  \node[lbl, anchor=south west] at (1.95,1.12) {$\vec n$ out of the page};
  \node[callout, anchor=north] at (1.8,-0.75) {fingers with the arrows, thumb along $\vec n$};
\end{tikzpicture}
\end{minipage}\hfill
\begin{minipage}[c]{0.46\linewidth}
\centering
{\small
\[
\begin{bmatrix} +1 & +1 & -1 & -1 \end{bmatrix}
\begin{bmatrix} e_1 \\ e_2 \\ e_3 \\ e_4 \end{bmatrix}
= \oint_{\partial A} \vec E \cdot \mathrm{d}\vec l
\]}
{\footnotesize the four non-zeros of one row of $\Cmat$; every other column is zero.\par}
\end{minipage}
\caption{One face of a brick and the four edges that bound it, traversed anticlockwise as seen from the tip of
the normal. Faraday's law for this face is the scalar equation above: the four edge integrals added in the sense
of the arrows equal minus the rate of change of the flux through the face. Because each edge has a positive
direction fixed by the grid axes, two of the four run against their stored sense and enter with a minus sign,
and the row of $\Cmat$ beside the picture records that pattern. Reading the row against the edge vector
reproduces the circulation; stacking one row per face gives the matrix equation in the text.}
\label{fig:hiw:faraday}
\end{figure}

The three subsections above are the whole of the discretisation, and Figure~\ref{fig:hiw:sequence} puts them
into a single line. The entities form a chain -- vertices, edges, faces, cells -- and the operators run along it:
the gradient takes vertex values to differences along the edges, the curl takes edge integrals to face fluxes,
the divergence takes face fluxes to the net outflow from a cell. Every arrow in that chain is exact, in the sense
that its image is precisely the kernel of the next one; that is what the two identities assert, and each arrow is
a matrix whose entries are $+1$, $-1$ or $0$ and nothing else. That is the claim the rest of this report leans
on, and it is a claim about incidence rather than about resolution.

\begin{figure}[htbp]
\centering
\begin{tikzpicture}[font=\footnotesize]
  \node[callout] (nd) at (0,0) {nodes $\approx N$};
  \node[callout] (ed) at (3.3,0) {edges $\approx 3N$};
  \node[callout] (fd) at (6.6,0) {faces $\approx 3N$};
  \node[callout] (cd) at (9.9,0) {cells $N$};
  \draw[oparrow] (nd.east) -- (ed.west);
  \draw[oparrow] (ed.east) -- (fd.west);
  \draw[oparrow] (fd.east) -- (cd.west);
  \node[lbl, anchor=south] at (1.5,0.10) {$\Gmat$: grad};
  \node[lbl, anchor=south] at (4.8,0.10) {$\Cmat$: curl};
  \node[lbl, anchor=south] at (8.1,0.10) {$\Smat$: div};
  \draw[oparrow] ($(nd.south)+(0,-0.05)$) .. controls ($(nd.south)+(0,-1.1)$) and ($(fd.south)+(0,-1.1)$) .. ($(fd.south)+(0,-0.05)$);
  \node[lbl, anchor=north] at (3.3,-1.05) {$\Cmat\Gmat=\mathbf{0}$: the loop telescopes};
  \draw[oparrow] ($(ed.south)+(0,-0.05)$) .. controls ($(ed.south)+(0,-2.0)$) and ($(cd.south)+(0,-2.0)$) .. ($(cd.south)+(0,-0.05)$);
  \node[lbl, anchor=north] at (6.6,-1.95) {$\Smat\Cmat=\mathbf{0}$: nothing leaks out of a closed surface};
\end{tikzpicture}
\caption{The chain of entities and the three operators that run along it, each with its matrix. The two
compositions underneath vanish identically. Read the picture as a statement about incidence and nothing else:
the loop integral of a gradient telescopes edge by edge, and the net flux of a curl out of a closed surface is
zero, for every mesh, because of what bounds what -- not because the mesh is fine.}
\label{fig:hiw:sequence}
\end{figure}

\subsection{Marching in time}\label{sec:hiw:leapfrog}

The two matrix equations have a property that is easy to miss and is the reason this is a time-domain method at
all. Faraday's law gives the rate of change of the face fluxes from the edge voltages and from nothing else;
Ampere's law gives the rate of change of the edge voltages from the face fluxes and the currents and from nothing
else. Each equation needs only the other family. So the two can be advanced in alternation: from the voltages at
one instant we update the fluxes, and from the updated fluxes we update the voltages. Neither half requires the
solution of a linear system, because the material matrices are diagonal, so every update is a division by a
number per unknown rather than an inversion of a matrix.

This is the trade the whole method makes, and it is worth naming the other side of it. A frequency-domain solver
inverts a large sparse matrix once per frequency and then has the whole spectrum; marching in time pays nothing
per step but must be run until the transient is over. What the time-domain route buys is the two things an
engineer usually wants from a structure of this kind: the response at every frequency at once, from one run, and
a field that can be watched while it evolves.

Written out, with $\mathbf{h}$ the magnetic field on the dual edges, $\mathbf{j}$ the current, and $\Delta t$ the
step,
\begin{equation}
  \begin{aligned}
    \mathbf{h}^{n+1/2} &= \mathbf{h}^{n-1/2} - \frac{\Delta t}{\mu}\,\Cmat\,\mathbf{e}^{n},\\
    \mathbf{e}^{n+1} &= \mathbf{e}^{n} + \frac{\Delta t}{\eps}\left(\Cmat^{\!\top}\mathbf{h}^{n+1/2}
       - \mathbf{j}^{n+1/2}\right).
  \end{aligned}
  \label{eq:hiw-leapfrog}
\end{equation}
The two families are offset by half a step -- the $\mathbf{h}$ values sit at $n\pm 1/2$ and the $\mathbf{e}$
values at $n$ -- which is where the name leapfrog comes from, and no unknown is ever updated from itself. That
is what keeps the scheme explicit and the storage down to two vectors per family.

The order of accuracy is second in time and second in space: halve the step and the brick and the error falls by
a factor of four. That is a statement about the method, and it is the reason the time step and the mesh have to
be refined together if the accuracy is to mean anything.

The step is not free. It is limited by the time light takes to cross the smallest brick, and in three dimensions
it is the diagonal of the brick that matters,
\begin{equation}
  \Delta t < \frac{\Delta x}{c\sqrt{3}} .
  \label{eq:hiw-courant}
\end{equation}
This bound is sharp, and it is worth being precise about how sharp. Measured on the implementation here, a run at
$1.01$ times the limit was stable and a run at $1.02$ times the limit diverged; Figure~\ref{fig:hiw:courant}
marks the two. Sharpness is a physical statement, not a numerical curiosity: below the bound the scheme does not
merely fail to blow up, it passes energy around consistently, one brick at a time; above it the update amplifies
what it is given and the field grows without bound.

The cost then follows by arithmetic rather than by measurement. In three dimensions, halving the brick gives
eight times the unknowns and halves the time step, so sixteen times the work per unit of simulated time. That is
the honest reason a method of this kind is used for one-off transient problems and for structures of a few
wavelengths, and not inside optimisation loops.

\begin{figure}[htbp]
\centering
\begin{tikzpicture}[font=\footnotesize]
  \node[edot] (e1) at (2.0,1.0) {};
  \node[edot] (e2) at (4.4,1.0) {};
  \node[edot] (e3) at (6.8,1.0) {};
  \node[lbl, anchor=south] at (2.0,1.12) {$\mathbf{e}^{n-1}$};
  \node[lbl, anchor=south] at (4.4,1.12) {$\mathbf{e}^{n}$};
  \node[lbl, anchor=south] at (6.8,1.12) {$\mathbf{e}^{n+1}$};
  \node[bdot] (h1) at (3.2,0) {};
  \node[bdot] (h2) at (5.6,0) {};
  \node[lbl, anchor=north] at (3.2,-0.12) {$\mathbf{h}^{n-1/2}$};
  \node[lbl, anchor=north] at (5.6,-0.12) {$\mathbf{h}^{n+1/2}$};
  \draw[oparrow] (e1) -- (h1);
  \draw[oparrow] (h1) -- (e2);
  \draw[oparrow] (e2) -- (h2);
  \draw[oparrow] (h2) -- (e3);
  \node[lbl, anchor=east] at (1.35,1.0) {primal edges};
  \node[lbl, anchor=east] at (2.55,0) {dual edges};
  \draw[oparrow] (1.4,-1.05) -- (7.9,-1.05);
  \node[lbl, anchor=west] at (7.95,-1.05) {$t$};
  \node[lbl, anchor=west] at (1.4,-1.75) {$\mathbf{h}^{n+1/2}=\mathbf{h}^{n-1/2}-\frac{\Delta t}{\mu}\,\Cmat\,\mathbf{e}^{n}$};
  \node[lbl, anchor=west] at (1.4,-2.45) {$\mathbf{e}^{n+1}=\mathbf{e}^{n}+\frac{\Delta t}{\eps}\left(\Cmat^{\!\top}\mathbf{h}^{n+1/2}-\mathbf{j}^{n+1/2}\right)$};
\end{tikzpicture}
\caption{The two families on a time axis. The magnetic field sits at half-integer levels, the electric field at
integer ones, and each staircase arrow consumes only the other family: no arrow starts and ends in the same row.
Read the figure as the reason the update is explicit -- there is never a moment at which an unknown is asked for
its own future value.}
\label{fig:hiw:leapfrog}
\end{figure}

\begin{figure}[htbp]
\centering
\begin{tikzpicture}[font=\footnotesize]
\begin{axis}[
  width=0.85\linewidth, height=0.18\linewidth,
  xmin=0, xmax=1.25, ymin=0, ymax=1,
  hide y axis, axis y line=none, ytick=\empty,
  xlabel={$\Delta t\,/\,(\Delta x/(c\sqrt{3}))$},
  tick label style={font=\scriptsize}, label style={font=\small},
  axis line style={draw=black!70}, tick style={draw=black!70},
  enlargelimits=false, clip=false,
]
  \fill[softgrey] (axis cs:0,0) rectangle (axis cs:1,1);
  \node[lbl] at (axis cs:0.42,0.5) {stable};
  \node[lbl, anchor=west] at (axis cs:1.04,0.5) {divergent};
  \draw[ref line] (axis cs:1,0) -- (axis cs:1,1);
  \addplot[only marks, mark=*, mark size=2.2pt, draw=efield, fill=efield]
    coordinates {(1.01,0.25) (1.02,0.75)};
  \node[lbl, anchor=west, font=\scriptsize] at (axis cs:1.04,0.25) {stable};
  \node[lbl, anchor=west, font=\scriptsize] at (axis cs:1.04,0.75) {diverged};
\end{axis}
\end{tikzpicture}
\caption{The stability bound as a hard edge rather than a rule of thumb. Everything to the left of the dotted
line is stable and nothing to the right of it is, as the two marks show: the time step that was stable sat at
$1.01$ times the limit and the next one, at $1.02$ times, diverged. There is no gradual degradation to trade
against -- the scheme either passes energy around consistently or it amplifies.}
\label{fig:hiw:courant}
\end{figure}

\subsection{What the machine is allowed to get wrong}\label{sec:hiw:materials}

The material enters the discrete problem through exactly two maps, $\mathbf{d}=\Em\mathbf{e}$ and
$\mathbf{b}=\Mm\mathbf{h}$, and nothing else. For a uniform material filling the whole box both matrices are
diagonal, and the update stays a division per unknown. Their entries are not simply the material constant,
however. In the diagonal case each entry carries the permittivity of its brick multiplied by a geometric factor,
the ratio of a face area to an edge length -- the cross-section of the dual cell to its length. That factor is
why a coarse mesh is wrong by a finite amount and a fine mesh converges: the material map is where the mesh
resolution actually enters the answer, and it is the only place it does.

An interface is where the difficulty sits. Put a brick of substrate against a brick of air and the edges lying in
the shared face belong to both; the discrete problem has one number to give each of them and they must be given
an average. Which average is a modelling decision, not a consequence of Maxwell's equations, and the choice is
visible in the result: two rules give answers that differ by amounts which do not shrink as the mesh is refined
unless the rule is the right one. The discontinuity in Figure~\ref{fig:hiw:interface} is real; the single value
the grid hands to an edge straddling it is a decision we make on the reader's behalf and should declare.

\emph{Every other ingredient on these pages is exact topology, and this is the only place where the discrete
problem is not the continuum problem.} Refining the mesh converges to the answer implied by the averaging rule
one has chosen, which is not necessarily the answer to Maxwell's equations.

That is not a technicality to be filed away. It is the question the later sections keep having to ask about
boxes and boundaries: the microstrip work in Section~\ref{sec:microstrip} cannot report a number without saying
what it did at the conductor edge, and the reference box in Section~\ref{sec:msres} is an attempt to make the
consequence of that choice measurable rather than assumed.

\begin{figure}[htbp]
\centering
\begin{tikzpicture}[font=\footnotesize]
  \draw[brick] (0,0) rectangle (2,1.6);
  \node[lbl] at (1,0.8) {substrate, $\eps_r$};
  \draw[mesh] (2,0) rectangle (4,1.6);
  \node[lbl] at (3,0.8) {air};
  \draw[mesh] (2,1.6) rectangle (4,3.2);
  \node[lbl] at (3,2.4) {air};
  \fill[face] (1.92,0) rectangle (2.08,3.2);
  \draw[edof] (2,0.38) -- (2,1.22);
  \node[edot] at (2,0.8) {};
  \node[callout, anchor=south] at (2,3.55) {one number must be chosen here};
  \draw[oparrow] (2,3.42) -- (2,3.12);
  % inset: eps against z across the interface
  \draw[mesh] (5.3,0.5) -- (5.3,2.4);
  \draw[mesh] (5.3,0.5) -- (8.0,0.5);
  \node[lbl, anchor=east] at (5.2,2.35) {$\eps_r$};
  \node[lbl, anchor=west] at (8.05,0.42) {$z$};
  \draw[primal] (5.5,2.0) -- (6.3,2.0) -- (6.3,0.95) -- (7.8,0.95);
  \draw[edof] (5.95,1.5) -- (6.65,1.5);
  \node[lbl, anchor=west] at (6.75,1.45) {what the edge receives};
\end{tikzpicture}
\caption{A cross-section through three bricks with a substrate on the left and air on the right and above, and
the edge that lies in the interface picked out in red. The grid has exactly one number to hand that edge, drawn
as the dashed horizontal value in the inset, while the permittivity it is trying to represent jumps. Everything
that can be wrong in this method is in this picture; nothing in the topological half of it can be.}
\label{fig:hiw:interface}
\end{figure}

\subsection{Driving it: what a source is}\label{sec:hiw:source}

A source is nothing more than a current term. The update displayed in Section~\ref{sec:hiw:leapfrog} advances
the electric field from a curl of the magnetic field plus a current on the same edges, and the term
$\mathbf{j}^{n+1/2}$ in it is the whole of the mathematics. There is no excitation module, no injection
operator and no port object in the discrete problem: whatever launches a wave -- a feed, a line port, a plane
wave introduced through a total-field boundary -- arrives in the update as a current on a set of edges or, in
the dual update, as a magnetic current on a set of faces. The physics does not care which notation was used to
write that current down. The numbers do, because the two obvious ways of adding it behave differently towards
what comes back.

The first way is a hard source: overwrite the unknown on the driven edges with a prescribed value at every
step, so that the value on them is what was said and not what the neighbours computed. It is the shortest thing
to write in a test and it is a lie about the physics. An edge whose value is dictated cannot respond to the
field around it, and a tangential electric field held fixed at every instant is the boundary condition of a
perfect conductor by definition. The driven edge is a metal wall of our choosing: it radiates, which is why the
lie is easy to live with and why the run looks like a wave leaving a source, and it reflects everything that
returns to it. Use it to check a code against a closed-form answer -- the prescribed value is then a function
whose analytic field is already known -- and never to measure a device.

The second way is a soft source: add the current and leave the update alone. The value on the driven edge is
still computed from its neighbours; the current is one more term in the sum that produces it. The edge
therefore still obeys Maxwell, still feels the field that comes back to it, and is pushed back on by the wave
it launched. The consequence that matters is this: a soft source has an impedance, and it does not absorb what
comes back. Whatever travels back through the driven edge is reflected by it a second time, in phase or out of
it depending on that impedance, and re-radiated. A run driven by a soft source inside a closed test window
contains a standing wave by construction -- not because the solver is wrong, but because the source is in the
box with the wave.

Figure~\ref{fig:hiw:hardsoft} draws the two side by side. Neither behaviour is a defect to be repaired: it is
what the words ``driven by a source at this point'' mean. It does mean two things have to be declared wherever
a trace is shown, namely which of the two ways was used and where the source sat relative to every surface the
wave could return from. Section~\ref{sec:sp-definition} is where the same difficulty reappears in the language
of measurement: a scattering parameter is defined only when the incident wave is the only one arriving at the
port, and the terminations and reference planes are what make that true.

\begin{figure}[htbp]
\centering
\begin{tikzpicture}[font=\footnotesize]
  % ---------------- (a) hard source ----------------
  \draw[mesh thin] (0.05,0.2) -- (0.05,2.4);
  \draw[edof] (1.9,0.3) -- (1.9,2.3);
  \node[edot] at (1.9,1.3) {};
  \node[lbl, anchor=west] at (2.05,2.15) {$e=\mathrm{fixed}$};
  \draw[edof] (0.4,1.3) -- (1.78,1.3);
  \node[lbl, anchor=south] at (1.05,1.38) {prescribed};
  \draw[meshgrey, line width=0.7pt, -{Stealth[length=2.6pt]}] (1.8,0.8) -- (0.55,0.8);
  \node[lbl, anchor=north] at (1.18,0.68) {reflected};
  \draw[mesh, densely dotted] (2.45,1.75) .. controls (2.15,1.6) .. (1.95,1.32);
  \node[lbl, anchor=west] at (2.05,1.78) {nothing crosses};
  \node[lbl, anchor=north] at (1.55,-0.1) {(a) hard: the wall dictates};
  % ---------------- (b) soft source ----------------
  \draw[mesh thin] (5.75,0.2) -- (5.75,2.4);
  \fill[face] (6.85,1.0) rectangle (7.45,1.7);
  \draw[edof] (7.15,0.85) -- (7.15,1.95);
  \node[lbl, anchor=west] at (7.3,2.0) {$\mathbf{j}$};
  \draw[harrow] (7.6,1.95) -- (9.6,1.95);
  \node[lbl, anchor=south] at (8.55,2.02) {launched};
  \draw[meshgrey, line width=0.7pt, -{Stealth[length=2.6pt]}] (9.55,0.75) -- (5.85,0.75);
  \node[lbl, anchor=north] at (8.3,0.63) {returns, and passes through};
  \node[callout, align=center, anchor=north] at (7.8,-0.1)
    {the source has an impedance;\\it does not absorb what comes back};
  \node[lbl, anchor=north] at (7.15,-1.15) {(b) soft: the source participates};
\end{tikzpicture}
\caption{The two ways of putting current into the update, drawn on the same edge. In (a) the value is
prescribed, so the edge is a metal wall: the field behind it cannot reach the unknown and what returns is
reflected. In (b) the current is added and the edge computes its own value, so the charge on it obeys its own
update and the returning wave passes through the source and continues rather than bouncing off it.}\label{fig:hiw:hardsoft}
\end{figure}

\subsection{Walls: what a box does to a wave}\label{sec:hiw:walls}

Every calculation in this report happens in a box, and a box has walls. A wall is a modelling decision of
exactly the same standing as the material average of Section~\ref{sec:hiw:materials}: the physics that comes
out is the physics of the box that was chosen, and a box is not free space. The walls are not a wrapper around
the real problem, added for the convenience of the solver; they are part of the problem, and a number taken
inside them is a number about the box as much as about the device.

Three kinds of wall are used, and each is worth stating twice -- once as what it does to the fields at the
wall, once as what it does to a wave arriving normally.

A perfect electric conductor sets the tangential electric field to zero. Nothing penetrates it, the surface
carries a current sheet, and the wave is reflected entirely with coefficient $\Gamma=-1$: the electric field of
the reflected wave points opposite to that of the incident one. A box lined with these is a cavity, so its
spectrum is discrete rather than continuous. On the grid the wall is the cheapest object there is: set the
tangential edges that lie on it to zero and never update them.

A perfect magnetic conductor sets the tangential magnetic field to zero instead. It is an open circuit rather
than a short, $\Gamma=+1$, and the electric field comes back in phase. No material is a magnetic wall at
microwave frequencies, and precisely that is what makes one useful. A structure symmetric about a plane can be
cut in half and the plane replaced by a wall -- which of the two walls it becomes depends only on whether the
drive is symmetric or antisymmetric about that plane. The full answer is then the half answer continued
across, at half the cost, and the wall is exact rather than approximate, because it encodes a symmetry of the
problem and not a property of any material.

A perfectly matched layer is not a boundary condition at all. It is a region of artificial loss, ideally
reflectionless, made by continuing the coordinates into the complex plane -- a complex coordinate stretch. The
reason it works is that a wave travelling in a coordinate with an imaginary part is attenuated while its
wavelength is unchanged, and because the stretch carries no frequency of its own the match is not tuned to a
band: it holds for whatever arrives, and the reflection does not grow as the layer is approached from any
angle. What the code implements is a conductivity profile whose strength ramps up over a few cells, with the
electric and magnetic losses held in the ratio that preserves the impedance of the medium; that ratio, not the
shape of the profile, is what does the matching.

A layer is only as good as its profile and its thickness, and both are choices rather than consequences. The
ramp in this work spans ten cells, which sits at the thin end of what is usual -- practice is commonly ten to
twenty -- and the thickness of the layer has never been swept here. That sentence is a statement about what has
not been measured, not about something known to be wrong, and it is the honest reason a wall reflection is
worth looking for in every trace.

Which is the practical point of this subsection. A wall reflection does not look like a phase error. It looks
like a standing wave: a ripple in the amplitude, with a period set by a length in the box rather than by
anything in the device, and it moves if the wall moves. A phase offset, by contrast, is a property of a
reference plane and travels with it. The diagnostic therefore needs no model at all -- if a curve ripples at a
period that matches a distance to a wall, the wall is talking, and the answer is a thicker layer or a farther
wall rather than a better description of the device. Figure~\ref{fig:hiw:walls} sets the three walls beside
one another; only the third returns nothing, and only the third is an approximation.

The formal treatment of the absorbing layer, its profile and the reflection it is measured to achieve, is in
Section~\ref{sec:res-pml}; the one-way condition it is compared against throughout is in
Section~\ref{sec:res-abc}.

\begin{figure}[htbp]
\centering
\begin{tikzpicture}[font=\footnotesize]
  % ---------------- PEC ----------------
  \draw[edof] (0.25,1.5) -- (2.15,1.5);
  \node[lbl, anchor=south] at (1.15,1.58) {incident};
  \fill[pattern=north east lines, pattern color=black!50] (2.25,0.15) rectangle (2.5,2.4);
  \draw[primal] (2.25,0.15) -- (2.25,2.4);
  \draw[edof] (2.1,0.85) -- (0.35,0.85);
  \node[lbl, anchor=north] at (1.2,0.72) {reflected, $\mathbf{E}$ reversed};
  \node[lbl, anchor=west] at (0.2,2.6) {$\hat{n}\times\mathbf{E}=0$};
  \node[lbl] at (1.4,-0.3) {$\Gamma=-1$};
  \node[lbl, anchor=south] at (1.4,3.05) {PEC};
  % ---------------- PMC ----------------
  \draw[edof] (4.05,1.5) -- (5.95,1.5);
  \node[lbl, anchor=south] at (4.95,1.58) {incident};
  \draw[dual] (6.05,0.15) -- (6.05,2.4);
  \node[lbl, anchor=west, align=left] at (6.12,2.1) {magnetic\\wall};
  \draw[edof] (5.9,0.85) -- (4.15,0.85);
  \node[lbl, anchor=north] at (5.0,0.72) {reflected, in phase};
  \node[lbl, anchor=west] at (4.0,2.6) {$\hat{n}\times\mathbf{H}=0$};
  \node[lbl] at (5.2,-0.3) {$\Gamma=+1$};
  \node[lbl, anchor=south] at (5.2,3.05) {PMC};
  % ---------------- PML ----------------
  \draw[edof] (7.95,1.5) -- (9.35,1.5);
  \node[lbl, anchor=south] at (8.55,1.58) {incident};
  \shade[left color=softgrey, right color=white] (9.45,0.15) rectangle (10.75,2.4);
  \draw[mesh] (9.45,0.15) rectangle (10.75,2.4);
  \draw[edof, dash pattern=on 2pt off 1.6pt] (9.5,1.5) -- (10.6,1.5);
  \node[lbl, anchor=west] at (7.85,2.6) {coordinate stretch};
  \node[lbl] at (9.5,-0.3) {$\Gamma\approx 0$};
  \node[lbl, anchor=south] at (9.5,3.05) {PML};
  \node[callout, align=center, anchor=north] at (10.1,-0.7) {ramp over\\a few cells};
  \draw[oparrow] (10.1,-0.68) -- (10.1,-0.05);
\end{tikzpicture}
\caption{Three walls and one wave arriving at each. The conductor reverses the electric field, the magnetic
wall returns it in phase, and the lossy ramp sends it into a coordinate that has an imaginary part and brings
nothing back. The first two are exact and visible in the answer; the third is the only one intended to be
invisible, and it is the only one that is an approximation.}\label{fig:hiw:walls}
\end{figure}


\subsection{Reading a number out: ports}\label{sec:hiw:ports}

A port is the place where it is decided what went in and what came out. There are two waves at a port, one
travelling in and one travelling out, and the measurement is the ratio of their amplitudes. Everything about
whether a number is quotable comes down to one question: can the port separate those two waves? If it cannot,
the ratio it reports is a mixture, and no amount of solver accuracy repairs a mixture.

The cheap construction is the lumped gap port. Leave a gap in the conductor, drive a current across it, read
the voltage as one line integral of the electric field across the gap -- a single edge -- and the current as
one loop integral of the magnetic field around the strip at that plane -- a single face loop. The impedance is
their ratio, $Z=V/I$, as drawn in Figure~\ref{fig:hiw:lumpedport}. This is a real measurement of the actual
fields at that plane, taken on the same unknowns the update produces; nothing is interpolated, fitted or
invented, and there is no hidden model between the field and the number.

It is not yet a scattering parameter, and the reason is intrinsic rather than numerical. $V$ and $I$ are two
different integrals over two different objects, an edge and a face loop, and their ratio equals the impedance
of a mode only if the field across the gap has exactly the transverse profile of that mode. A gap does not
excite a single profile. It excites a mixture, and the mixture evolves as it propagates: the components travel
at their own velocities and disperse, so the profile at the gap plane keeps changing, and it changes in
response to whatever comes back through it. $V/I$ therefore drifts as the profile evolves, and the drift is
contaminated by every other thing present at that plane, including whatever the box returns. That is a
statement about the port, not about the solver. The number is honest; it is simply not the number the device
is being asked about.

The construction that does give a scattering parameter is the modal port. Take the known transverse profile of
the mode, project the computed field at the port plane onto it -- a scalar product, an overlap integral -- then
inject only that amplitude going in and extract only that amplitude coming out. Now $V$ and $I$ are defined by
the projection rather than by one edge and one loop, and their ratio is a property of the mode rather than of
the gap. A modal port keeps one coefficient and discards the rest, and that discarding is what makes the
answer attributable to the device. The price is that the mode has to be known before the port can be built:
the profile calculation discussed with the power flow in Section~\ref{sec:sp-powerflow} is the thing that has
to have been done first, and in a guiding structure with an interface the profile is not something the
geometry alone will tell you.

One consequence has to be stated, because it governs how much a measured ratio can mean. Section~\ref{sec:hiw:source}
established that a soft source does not absorb the wave returning through it; inside a closed test window it
re-reflects that wave, so a forward and backward decomposition carried out in the same window sits on top of a
standing wave rather than on top of a travelling one. Any backward-to-forward ratio obtained that way is a
property of the source and the window as much as of the device between them, and the ambiguity does not
resolve with effort -- a finer mesh and a longer run both resolve the standing wave better without making it
smaller. What removes it is a port that separates the two waves by construction, together with walls that do
not return anything measurable (Section~\ref{sec:hiw:walls}).

The consequence for reading any number in this report is short. A scattering parameter is only as good as the
port that produced it, and ``the solver agreed with the analytic result'' has to mean the same port and the
same box on both sides of the comparison; otherwise the agreement is between two different questions that
happen to have similar answers.

\begin{figure}[htbp]
\centering
\begin{tikzpicture}[font=\footnotesize]
  % ---------------- the strip and its gap, seen from above ----------------
  \fill[brick] (0.3,1.25) rectangle (3.6,2.05);
  \fill[brick] (4.2,1.25) rectangle (7.5,2.05);
  \node[lbl, anchor=north] at (1.95,1.18) {strip};
  \node[lbl, anchor=north] at (5.85,1.18) {strip};
  \node[lbl, anchor=south] at (3.9,2.12) {gap};
  % one edge carries the voltage
  \draw[edof] (3.62,1.65) -- (4.18,1.65);
  \node[edot] at (3.9,1.65) {};
  \node[lbl, anchor=south] at (2.05,2.35) {$V=\int\mathbf{E}\cdot d\mathbf{l}$};
  \draw[oparrow] (2.55,2.28) -- (3.7,1.78);
  % one face loop carries the current
  \draw[dual] (3.9,1.65) ellipse [x radius=0.42, y radius=1.3];
  \draw[harrow] (4.3,2.6) -- (4.05,2.82);
  \node[lbl, anchor=west, align=left] at (4.35,2.85) {$I=\oint\mathbf{H}\cdot d\mathbf{l}$\\[-1pt]\scriptsize one face loop};
  \node[callout] at (3.9,0.55) {$Z=V/I$};
  \draw[oparrow] (3.9,0.72) -- (3.9,1.2);
  % ---------------- inset: the profile it assumes, and the one it has ----------------
  \draw[mesh] (9.1,0.55) -- (13.0,0.55);
  \draw[mesh] (9.1,0.55) -- (9.1,2.5);
  \draw[model line] plot[smooth] coordinates {(9.25,0.58) (9.95,2.3) (10.65,0.58)};
  \draw[efield, line width=0.7pt]
    plot[smooth] coordinates {(9.25,0.62) (9.55,1.55) (9.8,1.15) (10.05,2.1) (10.3,1.3) (10.65,0.66)};
  \node[lbl, anchor=west] at (10.7,2.15) {ideal mode profile};
  \node[lbl, anchor=west] at (10.7,1.0) {field at the gap};
  \node[lbl, anchor=east] at (9.02,2.42) {amplitude};
  \node[lbl, anchor=north] at (11.05,0.5) {transverse coordinate};
  \node[callout, align=center, anchor=north] at (11.05,-0.05)
    {these are not the same function,\\so $V/I$ is not the mode impedance};
\end{tikzpicture}
\caption{A lumped gap port. The voltage is one line integral across the gap and the current one loop integral
around the strip in the same plane, both taken on single grid objects, and the impedance is their ratio. The
inset is the assumption that ratio silently makes: that the profile at the gap is the profile of the mode. It
is not, which is why the number can be honest and still not be the one being asked for.}\label{fig:hiw:lumpedport}
\end{figure}

\begin{figure}[htbp]
\centering
\begin{tikzpicture}[font=\footnotesize]
  % ---------------- (1) whatever the field happens to be ----------------
  \draw[mesh] (1.1,0.35) -- (1.1,2.65);
  \draw[efield, line width=0.8pt]
    plot[smooth] coordinates {(1.15,0.5) (1.55,1.5) (1.85,1.05) (2.15,2.2) (2.45,1.2) (2.75,0.55)};
  \node[lbl, anchor=north, align=center] at (1.95,0.25) {field on the plane,\\whatever it is};
  % ---------------- (2) the basis and its coefficients ----------------
  \node[lbl] at (3.55,1.5) {$=$};
  \foreach \k/\y/\a in {1/2.05/0.55, 2/1.35/0.3, 3/0.65/0.15} {
    \draw[mesh thin] (4.0,\y) -- (6.1,\y);
    \draw[model line] plot[smooth] coordinates
      {(4.1,\y) (4.55,{\y+\a}) (5.0,\y) (5.45,{\y-\a*0.8}) (5.9,\y)};
    \node[lbl, anchor=west] at (6.15,\y) {$c_{\k}\,\phi_{\k}$};
    \node[lbl, anchor=east] at (3.92,\y) {\scriptsize mode \k};
  }
  % ---------------- (3) one coefficient kept ----------------
  \node[lbl] at (7.35,1.5) {$\longrightarrow$};
  \draw[mesh] (8.15,0.35) -- (8.15,2.65);
  \draw[model line] plot[smooth]
    coordinates {(8.2,0.5) (8.85,2.4) (9.5,0.5)};
  \draw[edof] (7.6,2.15) -- (8.05,2.15);
  \node[lbl, anchor=south] at (7.75,2.2) {inject};
  \draw[edof] (8.25,0.8) -- (7.6,0.8);
  \node[lbl, anchor=south] at (7.78,0.85) {extract};
  \node[lbl, anchor=north, align=center] at (8.85,0.2)
    {keep $c_{1}$ and discard the\\rest: the port now has\\a mode, and an impedance};
\end{tikzpicture}
\caption{The modal port as a projection. The field actually present on the plane is an arbitrary function; the
mode basis resolves it into coefficients, of which the port keeps exactly one. Injecting and extracting that
one coefficient -- and refusing to see the others -- is what turns a pair of integrals into a property of the
mode, and it is why the mode has to be known before the port can be built.}\label{fig:hiw:modalport}
\end{figure}

